\section{Related Work}

\paragraph{Neural operators on irregular geometries.}
Fourier neural operators learn resolution-transferable spectral mappings on
regular grids, while Geo-FNO extends them to general geometries through learned
coordinate deformations \citep{li2021fno,li2023geofno}. GINO couples integral
operators with a latent regular grid, and NUNO partitions non-uniform
observations to reduce interpolation and computational costs
\citep{li2023gino,liu2023nuno}. Transformer-based operators instead compress
large point sets into latent representations before global interaction, as in
GNOT, UPT, and Transolver \citep{hao2023gnot,alkin2024upt,wu2024transolver}.
Transolver++ and Transolver-3 further scale this design to million-scale and
industrial-scale geometries \citep{luo2025transolverpp,zhou2026transolver3}.
These methods establish strong compatibility with irregular inputs and changing
point counts. Our focus is complementary: we actively change the \emph{local}
sampling density at inference and ask whether the frozen model still represents
the same continuous operator.

\paragraph{Discretization-consistent operator learning.}
Extending finite neural architectures to function spaces requires discrete
reductions to converge to fixed continuous integrals as the discretization
changes \citep{berner2026functionspaces}. Monte Carlo convolution corrects
non-uniform point-cloud sampling with inverse-density weights, while QuadConv
evaluates continuous kernels using quadrature on irregular PDE data
\citep{hermosilla2018montecarlo,doherty2024quadconv}. Related formulations
construct discretization-invariant maps between neural fields and
quadrature-aware normalization statistics \citep{wang2022din,kim2026quadnorm}.
Nevertheless, resolution-independent parameters alone do not prevent errors
under a mismatch between training and evaluation discretizations
\citep{gao2025discretization}. We adopt the same function-space principle for
latent point aggregation, but use it to support a deliberate inference-time
change of sampling density rather than only passive resolution transfer.

\paragraph{Adaptive meshes and test-time refinement.}
Adaptive discretization concentrates resolution where numerical error is
expected to be large. Learned approaches predict mesh density, refinement, or
node relocation for numerical solvers
\citep{huang2021meshgeneration,obiols2023adarnet,rowbottom2024gadaptivity},
while MeshGraphNets integrates learned dynamics with mesh-based simulation
\citep{pfaff2021meshgraphnets}. Test-time computation has also been increased
through iterative denoising in PDE-Refiner, adaptive allocation or tokenization,
and neural-operator composition
\citep{lippe2023pderefiner,jenkins2026atca,
wang2026adaptivetokenization,serrano2026splitting}. These directions determine
where or how additional inference computation should be used. AdaSolver
addresses the complementary operator-level requirement: changing the spatial
sampling density of a frozen point-based neural operator while keeping its
reference continuous measure fixed.
